\documentclass{article}
\usepackage[UTF8]{ctex}
\usepackage{amsmath,amssymb,amsfonts}
\usepackage{enumitem}
\usepackage{geometry}
\geometry{a4paper, margin=1in}
% === 错题本排版优化（不参与题目标识） ===
\setlength{\parindent}{0pt}
\setlength{\parskip}{0.35em plus 0.10em minus 0.08em}
\setlength{\emergencystretch}{3em}
\setlist{itemsep=0.25em, topsep=0.40em, parsep=0pt}
\allowdisplaybreaks[2]
% === 错题本排版优化结束 ===
\begin{document}
\section*{吉林大学2024年高等代数真题}
\section*{68．吉林大学2024年高等代数与解析几何试卷}
三．设 \(A, B\) 均为实矩阵，证明：\(A^{\mathrm{T}} A+B^{\mathrm{T}} B\) 正定当且仅当 \(A X=0\) 与 \(B X=0\) 的公共解只有 0 解．

\begin{itemize}
\end{itemize}

\subsection{题目1}
原题边界“1．”：\(f, h\) 是数域 \(\Omega\) 上的首 1 多项式，\(f^{27} \mid h^{29}\) ，且 \(\operatorname{deg} h \leq 13\) ，证明：\(f \mid h\) ．

\subsection{题目2}
原题边界“2．”：设 \(V=\Omega^{n \times n}\) 是数域 \(\Omega\) 上所有 \(n\) 阶矩阵关于矩阵加法和数乘运算构成的向量空间， \(\mathcal{T}\) 满足
\[
\mathcal{T}(A)=E_{12} A E_{23}+E_{11} A E_{32}, \forall A \in \Omega^{n \times n} .
\]

\begin{itemize}
\item[（1）] 证明： \(\mathcal{T}\) 是 \(V\) 上的线性变换．
  \item[（2）] 求 \(\mathcal{T}(V)\) 的维数．
  \item[（3）] 给出 \(\mathcal{T}\) 的一个特征值。
\end{itemize}

四．设 \(V\) 为 \(n\) 维列向量空间，\(\sigma, \tau\) 均为 \(V\) 上的线性变换，\(\tau\) 为反称变换，且 \(\tau(V)=V\) 。证明：\(\sigma\) 为对称变换当且仅当 \(\tau \sigma \tau\) 为对称变换．

五 ．原题第（1）问：设 \(n\ge 4\)，并约定 \(E_{ij}\) 为 \(n\) 阶矩阵单位。证明：\(E_{12}+E_{34}\) 与 \(E_{1, n-1}+E_{1 n}+E_{2 n}\) 相似。\\

\subsection{题目3}
原题第（2）问：记 \(A=\left(\begin{array}{ccccc}a & 1 & & & \\ & a & \ddots & & \\ & & \ddots & 1 & \\ & & & a & 1 \\ & & & & b\end{array}\right)\) ，求 \(A\) 的 Jordan 标准型，其中 \(a \neq b\) ．

六．证明：\((\beta, \gamma, \delta) \alpha-(\delta, \alpha, \gamma) \beta+(\alpha, \beta, \delta) \gamma-(\alpha, \beta, \gamma) \delta=0\) ．

七．直线 \(\frac{x-x_{0}}{k}=\frac{y-y_{0}}{m}=\frac{z-z_{0}}{n}\) 在 \(y z, x z, x y\) 平面的投影直线记作 \(l_{1}, l_{2}, l_{3}\) ，原点与三条投影直线的距离记作 \(d_{1}, d_{2}, d_{3}\) ，原点与直线 \(\frac{x-x_{0}}{k}=\frac{y-y_{0}}{m}=\frac{z-z_{0}}{n}\) 的距离记作 \(d\) ，证明：

\[
d_{1}^{2}+d_{2}^{2}+d_{3}^{2}-d^{2}=k^{2} d_{1}^{2}+m^{2} d_{2}^{2}+n^{2} d_{3}^{2}
\]

八．证明：方程 \(e^{2 x+y+z}=x^{2}+y^{2}-2 x-4 y+4\) 是一个旋转曲面，并写出该旋转曲面的一条轴线．
\end{document}
